\documentclass[11pt,a4paper]{article}
\usepackage{amsmath,amssymb,amsthm}
\usepackage[margin=2.5cm]{geometry}
\usepackage{hyperref}

\newtheorem{definition}{Definition}[section]
\newtheorem{theorem}[definition]{Theorem}
\newtheorem{proposition}[definition]{Proposition}
\newtheorem{lemma}[definition]{Lemma}
\newtheorem{corollary}[definition]{Corollary}
\newtheorem{conjecture}[definition]{Conjecture}
\newtheorem{remark}[definition]{Remark}

\title{Hyperseed Formalization:\\
  Toward a Truly Decentralized Digital Provenance Layer}
\author{ProtomegaTron (automated formalization)}
\date{2026-09-18}

\begin{document}
\maketitle

\begin{abstract}
We formalize the intellectual structure of Ben Goertzel's Substack article
``Toward a Truly Decentralized Digital Provenance Layer''
(2026-08-25) within the Hyperseed ontology. The article introduces the
OpenWater protocol for decentralized digital provenance. We model
cryptographic provenance chains as proper fiber transport (in contrast to
watermarks as lossy base-space signals), decentralized trust as a flat
connection on the trust bundle, and the liar's dividend as a provenance
vacuum that the most powerful actor fills.
\end{abstract}

\tableofcontents

%% =================================================================
\section{The provenance crisis as fiber detachment}
\label{sec:crisis}
%% =================================================================

\begin{definition}[Digital artifact bundle --- claims 1--2]
\label{def:artifact-bundle}
Let $\pi : E \to B$ be a fiber bundle where:
\begin{itemize}
\item $B$ = the space of digital artifacts (images, audio, video, text),
\item $F_b = \pi^{-1}(b)$ = the provenance fiber over artifact $b$,
  containing the complete causal history of $b$: which cameras, microphones,
  keyboards, or generative models produced it, and what transformations
  it has undergone.
\end{itemize}
\end{definition}

\begin{proposition}[Provenance detachment --- claim 1]
\label{prop:detachment}
Modern generative technology has decoupled $B$ from $F$: for any
artifact $b \in B$, there exist multiple manufacturing paths
$p_1, p_2, \ldots \in F_b$ that produce perceptually identical outputs.
The projection $\pi$ is no longer injective on perceptual equivalence
classes:
\[
  b_{\text{real}} \approx_{\text{percept}} b_{\text{fake}}
  \quad \text{but} \quad
  F_{b_{\text{real}}} \neq F_{b_{\text{fake}}}
\]
The fiber carries the distinction; the base space does not.
\end{proposition}

\begin{corollary}[Scale defeats intuition --- claim 2]
Human perceptual judgment operates on $B$ (the artifact itself), not on
$F$ (its provenance). As generative fidelity increases, the information
available in $B$ alone becomes insufficient to distinguish real from
synthetic. ``Squint at it'' is a base-space heuristic applied to a
fiber-level problem.
\end{corollary}

%% =================================================================
\section{The liar's dividend as provenance vacuum}
\label{sec:liar}
%% =================================================================

\begin{definition}[Liar's dividend --- claim 3]
\label{def:liar}
When no artifact can be reliably authenticated (i.e., the provenance
fiber is inaccessible to all parties), the most powerful actor in a
dispute can deny any inconvenient evidence by asserting it belongs to
$F_{\text{fake}}$ without the accuser being able to prove otherwise:
\[
  \forall b,\; \text{if } F_b \text{ is inaccessible}: \quad
  \text{Power}(a) > \text{Power}(a') \implies
  a \text{ controls the narrative about } b
\]
\end{definition}

\begin{proposition}[Asymmetric corrosion --- claim 3]
\label{prop:corrosion}
The liar's dividend is more corrosive than deepfakes themselves because:
\begin{enumerate}
\item Deepfakes add false positives to the evidence pool (fixable by
  better detection).
\item The liar's dividend removes true positives by making all evidence
  dismissable (not fixable without provenance infrastructure).
\end{enumerate}
The second effect is structurally harder to counter because it exploits
the \emph{absence} of fiber access, not the presence of a specific fake.
\end{proposition}

%% =================================================================
\section{AI detection as base-space arms race}
\label{sec:detection}
%% =================================================================

\begin{proposition}[Detection is structurally losing --- claims 5--8]
\label{prop:detection-losing}
AI-based detection operates entirely in $B$: it examines artifact features
(pixel statistics, token distributions, stylometric patterns) and attempts
to infer $F_b$ from $b$ alone. This is a base-space inference about
fiber data, which is:
\begin{enumerate}
\item \textbf{Fundamentally underdetermined}: multiple fibers project to
  the same base-space region as generative fidelity increases.
\item \textbf{Co-evolutionarily unstable}: generators learn to evade
  detectors by matching the base-space statistics of authentic artifacts.
\item \textbf{Asymmetrically costly}: the attacker can train against the
  detector itself, making attack cost $\leq$ defense cost.
\end{enumerate}
\end{proposition}

%% =================================================================
\section{Cryptographic provenance as proper fiber transport}
\label{sec:openwater}
%% =================================================================

\begin{definition}[OpenWater provenance chain --- claims 12--14]
\label{def:openwater}
The OpenWater protocol attaches provenance data directly to the fiber
via cryptographic primitives:
\begin{itemize}
\item A \emph{provenance chain} $C(b) = (s_1, s_2, \ldots, s_n)$ is a
  sequence of signed claims, where each $s_i$ asserts that party $p_i$
  observed or transformed artifact $b$ at time $t_i$.
\item Each signature $s_i$ is verifiable without a central authority.
\item The chain $C(b)$ is metadata attached to $b$, not a perturbation
  of $b$'s content --- it survives content transformations.
\end{itemize}
\end{definition}

\begin{theorem}[Fiber transport preservation --- claim 13]
\label{thm:transport}
Let $T : B \to B$ be a content transformation (compression, transcoding,
cropping, translation). The OpenWater provenance chain satisfies:
\[
  C(T(b)) = C(b) \cup \{s_{T}\}
\]
where $s_T$ is the signature of the party applying $T$. The provenance
fiber \emph{grows} through transformations rather than being destroyed,
because it is fiber data (metadata) rather than base-space data
(content perturbation).

Compare with statistical watermarking (Section~5 of the companion
formalization), where $\rho(T) \ll 1$ --- the base-space signal is
destroyed by the same transformations.
\end{theorem}

\begin{remark}[Signed claims, not truth --- claim 14]
OpenWater does not claim to determine truth. It provides a verifiable
record of who claimed what about an artifact and when. Trust is built
from the network structure of claims, not from any single authoritative
assertion. This is the Hyperseed approach to provenance: provenance is
path data in the fiber, not a binary label.
\end{remark}

%% =================================================================
\section{Decentralized trust as flat connection}
\label{sec:flat}
%% =================================================================

\begin{definition}[Decentralized trust bundle --- claims 17--18]
\label{def:flat}
Let $\pi_T : E_T \to B_T$ be the trust bundle over the network of
provenance participants, with fiber $F_T = \{\text{trust level in
participant } p\}$. In a decentralized network:
\begin{itemize}
\item No participant has a privileged position (no ``Ministry of Reality'').
\item Trust is computed from the graph structure of mutual attestations.
\item The connection on $E_T$ is flat: parallel-transporting trust around
  any closed loop returns the same trust level.
\end{itemize}
\end{definition}

\begin{proposition}[No institutional holonomy --- claim 18]
\label{prop:no-holonomy}
In a centralized provenance system, trust can accumulate non-trivially
around institutional loops (cf.\ the EA holonomy in the companion
formalization). In the OpenWater decentralized network, there is no
closed institutional loop to traverse:
\[
  \mathrm{Hol}(\gamma) = \mathrm{id} \quad \forall \gamma \in \pi_1(B_T)
\]
The connection is flat because no single entity controls enough of the
network to create a monodromy.
\end{proposition}

%% =================================================================
\section{d-calculus perspective}
\label{sec:d-calc}
%% =================================================================

\begin{remark}[Hash chains as discrete parallel transport]
In the d-calculus framework, each cryptographic signature in the
OpenWater chain is a discrete parallel transport step: it maps the
provenance fiber from one transformation stage to the next while
preserving the accumulated history. The chain is a discretization of
the continuous parallel transport that Hyperseed uses to model
provenance flow through inference networks.
\end{remark}

\begin{remark}[Provenance as adaptive primitive]
The primitive decomposition of ``authentic content'' should include:
\texttt{capture-device-signature}, \texttt{transformation-chain},
\texttt{attester-network}, \texttt{temporal-ordering}. The current
media ecosystem operates with at most one primitive (``looks real''),
which is why deepfakes succeed. OpenWater implements the full
primitive set as cryptographic infrastructure.
\end{remark}

\begin{remark}[Subordinate AI role --- claim 11]
AI operates within the OpenWater framework as a fiber-consistency
checker: given a provenance chain $C(b)$, AI can verify that the
claimed transformations are consistent with the observed artifact
features. This is a fiber-to-base consistency check, not a base-only
inference --- a fundamentally more tractable problem than blind
detection.
\end{remark}

\end{document}
